\begin{textsample}{POS Dim 4 – 2009\_c5 – Score -1.00 – 2009\_c5/enade\_2009\_tec\_des\_mod\_3927bb5ec2.txt}  \label{ex:f4_pos_262}
Capítulo 5 Distribuição dos Conceitos 5.1 Panorama nacional da distribuição dos conceitos A Tabela 5.1 apresenta a quantidade e o percentual de cursos de TECNOLOGIA EM DESIGN DE MODA participantes do ENADE/2009 que se situaram em cada uma das faixas de conceitos, distribuídos de acordo com a Região em que se localizam os cursos das instituições de educação superior, mantenedoras dos cursos responsáveis pela oferta.
O Gráfico 5.1 apresenta a mesma informação sob a forma de diagrama de caixa para uma melhor visualização da dispersão dos conceitos obtidos.
No diagrama de caixa, a caixa propriamente dita fica limitada pelos quartis.
A barra dentro da caixa em negrito representa a mediana.
Linhas verticais unem os quartis aos valores extremos.
Se estes valores extremos se configurarem como valores atípicos, estes aparecem como pontos isolados e as linhas terminam no maior ( ou menor ) valor ainda não considerado atípico.
Eventuais valores atípicos aparecem como pequenos círculos com a Categoria Administrativa correspondente.
Valores considerados muito atípicos foram representados por estrelas.
Um diagrama de caixa para poucos dados ou com dados muito concentrados pode ter alguns pontos coincidentes.
Por exemplo, se mais de metade dos cursos de uma Região receberem conceito 5, o valor extremo superior, o quartil superior e a mediana coincidirão.
Se o conjunto só contém um dado, este apareceria denotado como a barra da mediana.
Se o conjunto apresenta todas as notas iguais, a notação é a mesma, somente uma barra horizontal.
Observando -se os dados da Tabela 5.1, percebe -se que, em termos nacionais, um pouco mais de 1/3 dos cursos ficaram sem conceito : 35,6 \% do total, correspondendo a 26 cursos do universo de 73 cursos.
No universo dos cursos que obtiveram conceito ( 47 ), predominaram os cursos que receberam conceitos 3 e 4, correspondentes, respectivamente a 31,9 \% e a 36,2 \% do total nacional avaliado.
Dos cursos, 2 foram classificados com conceito 1.
Os conceitos 2 e 5 foram alocados a, respectivamente, 12,8 \% e 14,9 \%.
Considerando -se separadamente as cinco Regiões, observou -se que, da Região Norte, somente 1 curso participou, o qual recebeu conceito 4.
Nenhum curso da Região obteve nem o conceito mínimo nem o conceito máximo.
Na Região Nordeste, 9 cursos participaram, dos quais, 4 ( 44,4 \% ) ficaram sem conceito.
Dentre os com conceito, os conceitos modais foram 3 e 4, correspondendo a 40 \% dos avaliados regionais.
Nenhum curso desta Região obteve o conceito mínimo e 1 curso ( 20,0 \% ) obteve conceito máximo. e 1 ( 10,0 \% ) obteve conceito 3.
Na Região Sudeste, 32 cursos participaram dos quais, 12 ( 37,5 \% ) cursos ficaram sem conceito.
Dentre os com conceito, o grupo modal foi o conceito 4, equivalentes a 35,0 \% do total regional avaliado.
O conceito máximo e o mínimo foram alcançados por 2 cursos cada ( 10,0 \% \% ) e 3 cursos foram classificados no conceito 3 ( 15 \% ) Na Região Sul, 10 dos 22 cursos avaliados, equivalentes a 45,5 \% do total regional, ficaram sem conceito.
Os 12 cursos avaliados ficaram distribuídos uniformemente nos conceitos 2, 3, 4 e 5 ( 3 cursos cada, correspondendo a 15 \% do total regional avaliado ) Na Região Centro-Oeste, 9 cursos participaram, dos quais nenhum ficou sem conceito.
Dos que obtiveram conceito, os conceitos modais foram 3 e 4 ( 44,4 \% equivalente a 4 cursos cada ).
O conceito máximo foi alcançado por somente 1 curso ( 11,1 \% ) e nenhum curso obteve o conceito mínimo. 5.2 Conceitos por Categoria Administrativa e por Região A Tabela 5.2 apresenta a distribuição dos cursos participantes do ENADE/2009 de TECNOLOGIA EM DESIGN DE MODA, por Categoria Administrativa, de acordo com os conceitos por eles alcançados, segundo Região.
O Gráfico 5.2 apresenta a mesma informação sob a forma de diagrama de caixa para uma melhor visualização da dispersão dos conceitos obtidos.
Somente eventuais valores atípicos aparecem com o identificador da Região na qual o curso foi habilitado.
De acordo com as informações apresentadas, a grande maioria, 90,4 \%, é de cursos em instituições privadas, os demais cursos estavam distribuídos entre as instituições federais, estaduais e municipais, com, respectivamente 2, 2 e 3 cursos.
Observando -se os dados da Tabela 5.2, percebe -se que, no universo dos 26 cursos que não obtiveram conceito, 1 era de instituição federal, 1 de instituição estadual, 2 de instituições municipais e os demais 22 de instituições privadas.
Na Região Norte, 1 curso em instituição privada não foi avaliado.
Na Região Nordeste dos 9 cursos sem conceito todos eram de instituições privadas.
Na Região Sudeste, os 32 cursos não avaliados eram também todos de instituições privadas.
Na Região Sul, dos 22 sem conceito, 1 era em instituição federal, 1 de instituição estadual, 1 de instituição municipal e os demais 17 de instituições privadas.
Na Região Centro-Oeste os 9 cursos sem conceito, 1 era em instituição federal, 1 de instituição estadual, e 7 eram, de instituições privadas.
A análise a seguir trata apenas da distribuição dos conceitos das 47 instituições que foram avaliadas, sendo que 1 em instituição federal, 1 em instituição estadual, 1 em instituição municipal e 44 em instituições privadas.
O único curso em instituição federal recebeu o conceito máximo, conceito 5.
Os únicos estadual e municipal receberam conceito 3.
Para os privados, o conceito modal para esta categoria foi 4, com 38,6 \% dos cursos avaliados.
O conceito mínimo foi alocado a 2 cursos.
O conceito máximo foi alcançado por 13,6 \% dos cursos em instituições privadas com conceito e 29,5 \% receberam conceito 3.
Considerando -se a análise por Região observou -se que em todas as Regiões brasileiras, houve representação de cursos em instituições privadas.
Na Região Norte, o único curso com conceito, 4, estava em instituição privada.
Obviamente nenhum curso recebeu nem o conceito mínimo nem o máximo.
Na Região Nordeste, apenas 5 cursos foram avaliados e todos em instituições privadas.
A distribuição foi : 2 com conceito 3, 2 com conceito 4 e 1 com máximo ( 5 ).
Nenhum recebeu o conceito mínimo.
Na Região Sudeste, 20 cursos e todos em instituições privadas receberam conceito : 2 com conceito mínimo, 3 com conceito 2, 6 com conceito 3, 7 com conceito 4 e 2 com conceito máximo ( 5 ).
Na Região Sul 12 cursos receberam conceito, apenas 1 em instituição municipal e o demais em instituições privadas.
O único municipal recebeu conceito 3.
Dos 11 cursos em instituições privadas, nenhum recebeu conceito mínimo, 3 receberam o conceito 2, 2 receberam conceito 3, 3 o conceito 4 e 3 o conceito máximo ( 5 ).
Na Região Centro-Oeste, 9 cursos receberam conceito, apenas 1 em instituição federal e 1 em instituição estadual, sendo que o demais 7 em instituições privadas.
O único federal recebeu conceito máximo e o estadual conceito 3.
Dos 7 cursos em instituições privadas, 3 receberam conceito 3 e 4 o conceito 4.
Nenhum recebeu nem conceito mínimo, nem máximo. 5.3 Conceitos por Organização Acadêmica e por Região Na Tabela 5.3 encontra -se a distribuição dos conceitos atribuídos aos cursos participantes do ENADE/2009 de TECNOLOGIA EM DESIGN DE MODA, por Organização Acadêmica, segundo as Regiões.
O Gráfico 5.3 apresenta a mesma informação sob a forma de diagrama de caixa para uma melhor visualização da dispersão dos conceitos obtidos.
Somente eventuais valores atípicos apareceram com o identificador da Categoria Administrativa.
Os cursos desta área estão principalmente alocados em Universidades ( 42,5 \% ) depois igualmente em Faculdades e Centros Universitários ( 28,8 \% cada ).
Nas Universidades, o conceito modal foi 3 ( 1/3 dos avaliados ), nos Centros Universitários o conceito modal ficou dividido entre 3 e 4 ( 41,7 \% dos avaliados ).
Nas faculdades, o conceito modal foi 4 com exatamente metade dos cursos avaliados neste conceito.
Os 26 cursos sem conceitos estão distribuídos em Universidades ( 10 ), em Centros Universitários ( 9 ) e em Faculdades ( 7 ).
Na Região Norte, o único curso avaliado pertencia a um Centro Universitário e recebeu conceito 4.
Nenhum curso ficou sem conceito na Região Norte.
Na Região Nordeste dos 4 cursos sem conceito, 2 estavam alocados em Centros Universitários e 2 em Faculdades.
Na Região Sudeste, dos 22 cursos não avaliados, 15 estavam alocados em Universidades, 3 em Centros Universitários e 4 em Faculdades.
Na Região Sul, dos 10 sem conceito, 7 estavam alocados em Universidades, 1 em Centro Universitário e 2 em Faculdades.
Na Região Centro-Oeste não houve curso sem conceito.
A análise a seguir trata apenas da distribuição dos conceitos dos 47 cursos que foram avaliados, sendo que 21 em Universidades, 12 em Centros Universitários e 14 em Faculdades.
Considerando -se separadamente as cinco Regiões, observa -se que, na Região Norte, como já comentado, o único curso pertence a um Centro Universitário e recebeu conceito 4.
A Região Nordeste apresentou 5 cursos na área de TECNOLOGIA EM DESIGN DE MODA com conceito : 1 em Universidade com conceito 3, 1 em Centro Universitário com conceito 4 e 3 em faculdades com conceitos 3, 4 e 5.
A Região Sudeste apresentou 20 cursos na área de TECNOLOGIA EM DESIGN DE MODA com conceito, sendo 7 em Universidades, 6 em Centros Universitários e 7 em Faculdades.
A categoria modal foi o conceito 2 nas Universidades ( 3 cursos ), 3 nos Centros Universitários ( 3 cursos ) e 4 em Faculdades ( 4 cursos ).
A Região Sul apresentou 12 cursos avaliados na área de TECNOLOGIA EM DESIGN DE MODA sendo 8 em Universidades e 2 cada em Centros Universitários e Faculdades.
O conceito modal nas Universidades foi o conceito máximo, 5 ( 3 cursos ).
Os dois cursos avaliados em Centros Universitários receberam conceito 3 e 4.
Os dois cursos avaliados em Faculdades receberam conceito 2.
Na Região Centro-Oeste, dos 9 cursos avaliados na área de TECNOLOGIA EM DESIGN DE MODA, 5 estavam alocados em Universidades e 2 cada em Centros Universitários e Faculdades.
O conceito modal para as Universidades foi 3 ( 3 cursos ).
Os dois cursos em Centros Universitários receberam conceito 3 e 4 e os dois cursos em Faculdades receberam conceito 4.

% matched lemmas: 
\end{textsample}
